% LiveGraph: ACL-style system demonstration paper (4 pages + references).
% Build: see README.md in this directory (needs acl.sty / acl_natbib.bst from acl-org/acl-style-files).
\documentclass[11pt]{article}
\usepackage[review]{acl}

\usepackage[T1]{fontenc}
\usepackage[utf8]{inputenc}
\usepackage{microtype}
\usepackage{inconsolata}
\usepackage{booktabs}
\usepackage{graphicx}
\usepackage{xcolor}
\usepackage{amsmath}

\newcommand{\todo}[1]{\textcolor{red}{[TODO: #1]}}
\newcommand{\sys}{\textsc{LiveGraph}}
\newcommand{\results}{../results}
% Figures: `report` writes SVG; convert to PDF for LaTeX (see README.md), else a placeholder box is shown.
\newcommand{\resultfig}[2]{%
  \IfFileExists{\results/figures/#1.pdf}%
    {\includegraphics[width=#2]{\results/figures/#1.pdf}}%
    {\fbox{\parbox{0.9\linewidth}{\centering\small Figure pending: \texttt{#1}}}}}
\newcommand{\resulttable}[1]{%
  \IfFileExists{\results/tables/#1.tex}{\input{\results/tables/#1.tex}}{\todo{table #1 pending}}}

\title{\sys: A Live, Steerable Multi-Agent Research Canvas\\with Incremental Recompute}

\author{Anonymous submission}

\begin{document}
\maketitle

\begin{abstract}
Deep-research agents turn a question into a cited report, but they run as opaque one-shot pipelines:
users cannot see which evidence a claim rests on, how confident the system is in it, or fix one step
without paying for, and waiting for, the whole run again. We present \sys{}, an open-source research
canvas in which a planner agent builds the research workflow as a graph on screen, every step streams
live, and the user can steer, edit, rerun, restore or veto any step, or distrust a source. Claims are
contested by a prosecutor and a defender and scored by a judge with per-claim confidence. Every node
output is an immutable version keyed by a hash of its inputs, so a change recomputes only the affected
steps, and an early cutoff stops propagation when a recomputed output is unchanged. On \todo{N} research
questions, \sys{} \todo{factual precision result} and updates a brief after a change at
\todo{X}\% of the cost of starting over. Code, evaluation harness and data are released.
\end{abstract}

\section{Introduction}
Systems that research a question on the web and write a cited report (``deep research'') are now
widely used. Their interface, however, is a chat message in and a long document out. The intermediate
work (which sub-questions were explored, what evidence was found, which claims were contested) is hidden
or shown only as a scrolling log. When a reader spots a problem, such as a claim resting on an unreliable
source, their only recourse is to re-prompt, which reruns the whole pipeline, costs as much as the first
run, and may change parts of the report that were fine.

We argue that a research report should behave like a build: a graph of steps with explicit inputs, where a
change to one input rebuilds only what depends on it. \sys{} applies this idea to multi-agent research.
Its contributions are:
\begin{itemize}
  \item \textbf{A live, inspectable workflow graph.} A planner agent creates the graph through tool calls,
  so nodes appear on the canvas as it plans; each node streams its output and its tool activity, and a
  timeline scrubber replays the full run from a persisted event log (\S\ref{sec:system}).
  \item \textbf{Adversarial verification with calibrated confidence.} Extracted claims are attacked by a
  prosecutor and defended by a defender; a judge, who can ask either side for rebuttals, assigns each claim
  a confidence or strikes it, and the editor hedges or drops claims accordingly.
  \item \textbf{Incremental recompute for agent workflows.} Node outputs are immutable versions keyed by an
  input hash; steering, editing or distrusting a source invalidates only downstream nodes, and an early
  cutoff on output hashes stops recomputation as soon as an output is unchanged.
  \item \textbf{An open evaluation harness} (\texttt{lgeval}) measuring factual precision, citation quality,
  calibration and recompute savings, with a pre-registered protocol (\S\ref{sec:eval}).
\end{itemize}

\section{System}
\label{sec:system}
\begin{figure}[t]
  \centering
  \fbox{\parbox{0.95\linewidth}{\centering\small \todo{Screenshot: canvas mid-run with claims panel}}}
  \caption{The \sys{} canvas during a Living Brief run. \todo{annotate}}
  \label{fig:screenshot}
\end{figure}

\paragraph{Living Brief workflow.} The user enters a goal and, optionally, uploads documents. A
\emph{planner} agent decomposes the goal into 3--5 sub-questions and builds the graph with
\texttt{add\_node}/\texttt{add\_edge} tools; each call is persisted and broadcast immediately, so the graph
grows on screen. \emph{Researcher} agents search the uploaded corpus and the web, read sources and report
findings tied to chunk ids; a \emph{critic} checks each researcher and can add a follow-up research node
while the run is in progress. A \emph{claims} step extracts atomic claims, each with supporting chunks. A
\emph{prosecutor} and a \emph{defender} argue against and for each claim with evidence, and a \emph{judge}
reads both, may request targeted rebuttals, and outputs a confidence in $[0,1]$ per claim or strikes it.
Finally an \emph{editor} writes the brief, citing claims inline, excluding struck claims and hedging
low-confidence ones. Source text reaches models only as delimited data, and output validators reject
citations to chunks the node never saw.

\paragraph{Versioned nodes, input hashes and early cutoff.} Every node output is stored as an immutable
version. Before a node runs, the engine computes an input hash over its type, configuration, the content
hashes of its inputs, applicable steering text and, for corpus-reading nodes, a fingerprint of the
documents in scope and their trust levels. If a version with that hash exists, it is reused without a model
call. Each version also has an output hash over its content and cited evidence. When a recomputed node
produces an output with the same hash as before, its descendants are not invalidated; those already marked
stale re-hash, find their existing version and finish without running. This is the early-cutoff
optimisation of build systems \citep{mokhov2018buildsystems} applied to LLM agents: an identical rerun
stops at the rerun node. Web results are cached per run, so recomputation does not re-query the web
unless the user asks for fresh sources.

\paragraph{Steering and editing.} Any node can be steered with a natural-language instruction (a running
node restarts with it; a finished node produces a new version), edited directly, rerun, restored to an
earlier version or vetoed. Distrusting or pinning a document changes the corpus fingerprint of the nodes
that read it, so exactly those nodes and their dependents recompute. An impact preview shows which nodes an
action may touch before the user commits to it. Downstream agents can also request work from upstream
nodes (for example, a judge asking a researcher to look again), which flows through the same mechanism.

\paragraph{Multi-viewer and replay.} Every state change is an event written to Postgres before it is
broadcast over WebSockets, so any number of viewers see the same run live, and the state at any point in
time is a fold over the event log, which powers the timeline scrubber and version diffs.

\paragraph{Implementation.} The backend is FastAPI with Pydantic AI agents over OpenRouter
(any OpenAI-compatible model per role; in all experiments every agent runs on \texttt{deepseek-v4.1-flash}),
Postgres with pgvector for the corpus; the frontend is Next.js with React Flow. A typical run costs about
\$0.01 \todo{update with measured pilot cost}. \todo{link to code and demo video}

\section{Evaluation}
\label{sec:eval}
We test three claims (protocol in the repository, \texttt{docs/EVAL.md}): (A) briefs are more factual and
better cited than ablations and a single-agent baseline with the same models and tools; (B) judge
confidence is calibrated; (C) updates after a change cost a fraction of starting over without quality
loss. We use \todo{N} custom research questions and \todo{N} questions from FRAMES
\citep{krishna2024frames}, three repeats per system. Factual precision follows FActScore and SAFE
\citep{min2023factscore,wei2024longform}, citation precision/recall follows ALCE \citep{gao2023alce}, and
short answers are graded as in SimpleQA \citep{wei2024simpleqa}. Scores come from an LLM judge of a
different model family than \sys{}'s models, with position-swapped pairwise comparisons
\citep{zheng2023judging} and a human spot-check of \todo{X}\% of decisions (Cohen's $\kappa$ = \todo{}).
CIs are 95\% bootstrap intervals over questions.

\resulttable{quality}

\paragraph{Quality (A).} \todo{Table~\ref{tab:quality}: summarise full vs.\ ablations and baseline.}
\resulttable{pairwise}

\paragraph{Calibration (B).} \todo{Brier, ECE, AUROC from Table~\ref{tab:calibration}; Figure~\ref{fig:reliability}.}
\resulttable{calibration}
\begin{figure}[t]
  \centering
  \resultfig{reliability}{0.85\linewidth}
  \caption{Reliability diagram of judge confidence (marker area $\propto$ claims per bin).}
  \label{fig:reliability}
\end{figure}

\paragraph{Incremental recompute (C).} \todo{Table~\ref{tab:recompute}: cost/time \% per mutation, nodes reused.}
\resulttable{recompute}
\begin{figure}[t]
  \centering
  \resultfig{recompute_cost}{0.85\linewidth}
  \caption{Mean cost of an update vs.\ starting over, per mutation.}
  \label{fig:recompute}
\end{figure}

\paragraph{Cost of the study.} The full evaluation cost \todo{\$X} (Table~\ref{tab:study_cost}).
\resulttable{study_cost}

\section{Related Work}
\paragraph{Deep research and report generation.} STORM \citep{shao2024storm} generates Wikipedia-like
articles by simulating perspective-guided conversations to collect sources before outlining and writing.
Commercial deep-research products follow the same one-shot pattern. Benchmarks such as FRAMES
\citep{krishna2024frames} and BrowseComp \citep{wei2025browsecomp} evaluate retrieval and browsing on
questions with short answers. \sys{} focuses instead on the interaction model: making each step visible,
steerable and cheap to revise.

\paragraph{Debate, critique and self-refinement.} Debate was proposed as a way for a weaker judge to
supervise stronger agents \citep{irving2018debate}, and multi-agent debate improves factuality and
reasoning \citep{du2024debate}. Self-Refine \citep{madaan2023selfrefine} and Reflexion
\citep{shinn2023reflexion} improve outputs through self-feedback. \sys{} uses role-specialised
prosecutor and defender agents with a judge that outputs per-claim confidence, and exposes the debate in
the interface rather than only using its outcome.

\paragraph{Evaluating factuality and citations.} FActScore \citep{min2023factscore} and SAFE
\citep{wei2024longform} decompose long-form text into atomic facts; ALCE \citep{gao2023alce} measures
citation precision and recall; LLM-as-a-judge \citep{zheng2023judging} enables pairwise comparison but
exhibits position and self-preference biases \citep{panickssery2024selfpref}. We measure calibration with
Brier score and ECE \citep{brier1950,naeini2015ece,guo2017calibration}.

\paragraph{Incremental computation.} Adapton \citep{hammer2014adapton} provides demand-driven
incremental computation; build systems such as those formalised by \citet{mokhov2018buildsystems} rebuild
only out-of-date targets and use early cutoff when a rebuilt target is unchanged. We apply these ideas to
a graph whose steps are non-deterministic LLM agents, where reuse is keyed on input hashes and cutoff on
output hashes.

\section*{Limitations}
LLM outputs are non-deterministic, so ``unchanged'' is detected by exact output hashes; a semantically
equivalent but reworded output still propagates. Our quality metrics rely on LLM judges, which we mitigate
but cannot eliminate with a cross-family judge and a human spot-check. Factual precision is checked against
retrieved evidence, so errors present in sources are not caught. The evaluation uses \todo{N} questions in
English with one model configuration, and web search results drift over time. The user study is small
(\todo{n} participants) and formative. Commercial deep-research systems change without notice and are
compared only descriptively.

\section*{Ethics Statement}
\sys{} retrieves and summarises web content; briefs can contain errors, and confidence scores can be
miscalibrated, so the interface shows evidence and confidence rather than presenting outputs as fact.
Retrieved content is treated as data, not instructions, to reduce prompt-injection risk, but this is not a
guarantee. The user study collected informed consent, stored only pseudonymous interaction logs, and
\todo{IRB/ethics approval status}. We report the monetary cost of all experiments.

\bibliography{refs}

\end{document}
